\scanpage{23}
% sheet: reused_typed
% hand: typed
% archive sheet number: 22
% Typescript rotated 180 degrees relative to the foliation.
% @ text / head in reading orientation
\mspage{- 53 -}
% @ text / first paragraph, continuing from another leaf
holomorphe dans $P$ (comme fonction réciproque de $e^z$), on a $\handwritten{\int_\gamma}\frac{dz}{z}=\log(\gamma(b))-\log(\gamma(a))$, et par suite $\handwritten{\unsure{\mathfrak I|\mathscr I}}(\handwritten{\int_\gamma}\frac{dz}{z})=\mathrm{Am}(\gamma(b))--\mathrm{Am}(\gamma(a))$, en désignant par $\mathrm{Am}(z)$ la mesure de l'amplitude de $z$ appartenant à l'intervalle $\handwritten{]}-\pi,+\pi\handwritten{[}$ . Par abus de langage, on dit encore que $\handwritten{\mathfrak I}(\handwritten{\int_\gamma}\frac{dz}{z})$ pour une route \uline{quelconque} $\gamma$ contenue dans $\handwritten{\complement\{}0\handwritten{\}}$ , est la \uline{variation de l'amplitude de} $z$ \uline{le long de} $\gamma$ . Si $\gamma$ est un circuit d'indice $n$ par rapport à 0 ; la "variation de l'amplitude de z " le long de $\gamma$ est donc $2n\pi$ .

% @ text / middle, paths and deformations
Les notions précédentes relatives aux chemins \uline{carrossables} peuvent s'étendre aux chemins \uline{quelconques} . Soit $\gamma$ un chemin achevé quelconque défini dans un intervalle $I=\handwritten{[}\alpha,\beta\handwritten{]}$, à valeurs dans $\handwritten{\mathbb C}$ , et soit $a$ un point de $\handwritten{\mathbb C}$ n'appartenant pas à $\gamma(I)$ . Il existe une déformation de $\gamma$ en un chemin carrossable $\gamma'$ \uline{dans} $\handwritten{\complement\{}a\handwritten{\}}$ ($\S\,2$, prop.1) conservant les extrémités de $\gamma$ ; en outre, si $\gamma'$ et $\gamma''$ sont deux chemins carrossables déformés de $\gamma$ dans $\handwritten{\complement\{}a\handwritten{\}}$ par de telles déformations, $\gamma'$ et $\gamma''$ peuvent être déformés l'un en l'autre dans $\handwritten{\complement\{}a\handwritten{\}}$ par une déformation de même type ; en raison du cor. du th.1 du $\S\,2$ , la "variation de l'amplitude" de $z-a$ le long de $\gamma'$ et le long de $\gamma''$ sont égales. Cette variation ne dépend donc que du chemin $\gamma$ initial, et on l'appelle encore la \uline{variation de l'amplitude de} $z-a$ \uline{le long de} $\gamma$ .
% @ margin note / left of the four lines beginning peuvent, no letters
\margin{\handwritten{$\bigl(\ \big|$}}

% @ text / lower third, reversal and juxtaposition
Il résulte aussitôt de ces définitions que les variations de l'amplitude de $z-a$ le long de $\gamma$ et le long de l'opposé de $\gamma$ sont \uline{opposées} ; si $\gamma_1$ et $\gamma_2$ sont deux chemins dans $\handwritten{\complement\{}a\handwritten{\}}$ tels que la juxtaposition de ces chemins soit définie, la variation de l'amplitude de $z-a$ le long de $\gamma_1\handwritten{\vee}\gamma_2$ est la \uline{somme} des variations de l'amplitude de $z-a$ le long de $\gamma_1$ et le long de $\gamma_2$ .
